\bookchapter{Приложение C: Глоссарий технических терминов}

\begin{center}
\begin{tabularx}{\linewidth}{|X|X|X|}
\hline
\textbf{Термин} & \textbf{Что это на самом деле} & \textbf{Краткое объяснение} \tabularnewline
\hline
SoC / чип & Мозг и сердце телефона & Ключевой компонент, определяющий скорость телефона и плавность игр. \tabularnewline
RAM (оперативная память) & «Короткий рабочий стол» телефона & Чем больше «стол», тем больше приложений можно открыть одновременно и тем плавнее переключение. \tabularnewline
ROM (хранилище) & «Большой склад» телефона & Место для хранения фото, видео и приложений. 128 ГБ и 256 ГБ говорят именно о нём. \tabularnewline
OLED / LCD & Два типа материалов экрана & У OLED ярче цвета и ниже энергопотребление; LCD мягче для глаз и дешевле в производстве. \tabularnewline
NFC & Ближняя беспроводная связь & Позволяет телефону превращаться в транспортную карту, карту доступа, банковскую карту — функция «прикладывания». \tabularnewline
OTA-обновление & Беспроводное обновление системы & То самое «обновление системы», которое вы нажимаете в [Настройках]. \tabularnewline
Root & Права суперпользователя Android & После получения можно изменять системные файлы, но есть и риски. \tabularnewline
Блокировка BL & Блокировка Bootloader & Защитная блокировка от производителя; только после разблокировки можно прошивать систему. \tabularnewline
NPU & Нейросетевой процессор & Чип, специализированный для задач ИИ; в десятки раз быстрее CPU и с меньшим энергопотреблением. \tabularnewline
Локальный ИИ & ИИ, работающий на устройстве & ИИ запускается прямо на телефоне без подключения к сети; высокая скорость и хорошая конфиденциальность. \tabularnewline
Qi2 & Стандарт беспроводной зарядки нового поколения & Беспроводная зарядка с магнитным позиционированием, похожая на MagSafe; заряжать удобнее. \tabularnewline
eSIM & Встроенная SIM-карта & SIM без физической карты, прописывается прямо в телефон; для смены оператора не нужно менять карту. \tabularnewline
LTPO & Технология адаптивной частоты обновления & Частота обновления экрана подстраивается под контент (1–120 Гц); плавно и энергоэффективно. \tabularnewline
\hline
\end{tabularx}
\end{center}
